\begin{theorem}[Large-order terminal witness]
Let (V) and (E) be finite types with decidable equality, let
(G=(V,E)) be a loopless multigraph, and let (k\in\mathbb N).  Suppose
that (G) is \noderef{Critical},
[
 \noderef{ChromaticIndex}(G)=k+1,
 \qquad \noderef{MaximumDegree}(G)+1\leq k,
 \qquad 11\leq \noderef{TashkinovOrder}(G).
]
Then there exist an exact edge identity (e_0\in E), a partial
(k)-edge-coloring (\varphi), and a finite set (X\subseteq V) such
that (\varphi) is \noderef{IsProperEdgeColoring}, it leaves exactly
(e_0) uncolored in the sense of \noderef{IsOneEdgeUncolored}, every
endpoint of (e_0) belongs to (X), and (X) is both
\noderef{ColoredVertexSetElementary} and
\noderef{ColoredVertexSetStronglyClosed} with respect to (\varphi).
\end{theorem}
\begin{proof}
We first record an endpoint fact that will apply to every tree returned
below.  If (R) is a \noderef{TreeSequence} based at the exact edge
(e_0), write its first displayed edge as (e_1).  The definition gives
(e_1=e_0), and the attachment condition at the first positive edge index
says that the two ends of (e_1) are its new vertex (y_1) and a vertex
(y_i) with (i<1).  Necessarily (i=0).  Thus both exact endpoints of
(e_0) occur among the displayed vertices of (R), and hence, by
\noderef{TreeSequenceVertexSet},
[
                 v\in G.\operatorname{ends}(e_0)
                 \quad\Longrightarrow\quad v\in V(R).       \tag{1}
]

Apply \noderef{AlgorithmThreeOneInitialSeriesExists} to criticality, the
chromatic-index equality, and the degree bound.  It supplies an exact edge
(e_0), a partial coloring
(\varphi), and a certified zero-round series
(\mathcal S^0), together with a \noderef{GlobalMaximumTriple} based on
these data.  By that definition the underlying data form a
\noderef{KTriple}; hence (\varphi) is proper and leaves exactly (e_0)
uncolored.  The initial-series theorem also says that the unique coloring slot is
(\varphi), and that the terminal tree (T^0_1) has
[
       |V(T^0_1)|=\noderef{TashkinovOrder}(G).                 \tag{2}
]
It also says that (V(T^0_1)) is elementary and closed under
(\varphi).

If (V(T^0_1)) is strongly closed under (\varphi), take
(X=V(T^0_1)).  Properness, the exact-one-uncolored property,
elementarity, and strong closedness are the facts just obtained, while
(1), applied to (T^0_1), puts both ends of (e_0) in (X).  This is
the desired conclusion in this case.

Assume henceforth that (V(T^0_1)) is not strongly closed.  The exact
last coloring and tree indices of the zero-round series are respectively
its slots (0) and (1), so its terminal properness, unique-uncoloredness,
closedness, and failure of strong closedness meet all hypotheses of
\noderef{AlgorithmThreeOneStepExists}.  Apply that theorem at round zero.
It gives a certified one-round series (\mathcal S^1) whose data are the
exact \noderef{TashkinovSeriesDataAppend} of the zero-round data.  By
\noderef{TashkinovSeriesFinalETTExists}, the exact final tree (T^1_2),
the exact terminal coloring (\varphi^1_1), and this same dependent
series determine a positive ETT (\mathcal E^1) of rung exactly (1).

We verify \noderef{MaximumProperty} for (\mathcal E^1).  In the append,
the old terminal tree (T^0_1) is retained literally in the new tree slot
(T^1_1).  Therefore (2) gives
[
                         |V(T^1_1)|=t(G),                     \tag{3}
]
which is the first clause of the maximum property.  Its stagewise clause
ranges over indices (i) with (1\leq i<1), of which there are none.
Thus (\mathcal E^1) has the maximum property.

For (n\in\mathbb N), call (n) attainable if there exist an input
partial coloring (\rho), a tree-sequence (T), a generating coloring
(\sigma), and a positive ETT
[
                    \mathcal E:(G,k,e_0,\rho,T,\sigma)
]
such that the rung of (\mathcal E) is (n), the displayed tree (T)
is exactly the final tree of the dependent certified series of
(\mathcal E), and (\mathcal E) has the maximum property.  The ETT
(\mathcal E^1) proves that (1) is attainable.  Applying
\noderef{CertifiedSeriesRoundsBound} to its dependent one-round series
gives (1+2\leq |V|), and in particular (1<|V|).

Define the finite set
[
        A=\{m\in\mathbb N:m<|V|\text{ and }m\text{ is attainable}\}.
]
The preceding paragraph shows (1\in A), so (A) is nonempty.  Let
(n=\max A), and choose witnesses
((\rho_n,T_n,\sigma_n,\mathcal E_n)) for the attainability of (n).
In particular, (T_n) is the exact final tree of the certified series of
(\mathcal E_n), and (\mathcal E_n) has the maximum property.

Apply \noderef{PositiveWellDefinedness} to (\mathcal E_n), using this
exact-final equality, its maximum property, and the assumed degree and
large-order bounds.  Consider its two alternatives.

In the first alternative, there are a replacement input coloring
(\rho^*), a certified same-rung series (\mathcal S^*), and an ETT
(\mathcal E^*) displayed on exactly the final tree of (\mathcal S^*),
with generating coloring exactly its terminal coloring (\varphi^*).
The same alternative states that (\varphi^*) is proper and leaves exactly
(e_0) uncolored, that the final vertex set
[
 X=V\bigl(\mathcal S^*.T_{n+1}\bigr)
]
is strongly closed, and that (\mathcal E^*) has the maximum property.
Apply assertion (i) of \noderef{TechnicalTheoremLarge} to this exact ETT.
Its hypotheses are precisely that maximum property and the same degree and
large-order inequalities, and its conclusion is that (X) is elementary
under (\varphi^*).  Finally (1), applied to the exact final tree of
(\mathcal S^*), puts both endpoints of (e_0) in (X).  Hence
((e_0,\varphi^*,X)) has all five required properties.

It remains to rule out the second alternative.  That alternative supplies
a certified series (\mathcal S^+) with (n+1) noninitial iterations and
an exact-final ETT (\mathcal E^+).  More precisely, it states both that
the rung of (\mathcal E^+) is exactly (n+1) and that its dependent
series is exactly (\mathcal S^+); it also states that
(\mathcal E^+) has the maximum property.  Eliminating those two exact
dependent equalities identifies its displayed tree with the final tree of
(\mathcal S^+).  Consequently these data prove that (n+1) is
attainable.

Moreover, \noderef{CertifiedSeriesRoundsBound} applied to
(\mathcal S^+) gives
[
                         (n+1)+2\leq |V|,
]
so (n+1<|V|).  Thus (n+1\in A).  Since (n=\max A), this implies
(n+1\leq n), an arithmetic contradiction.  The extension alternative is
therefore impossible, and the witness obtained in the first alternative
completes the proof.
\end{proof}
